\scanpage{135}

lemme 1 et l'hyp. de récurrence, OK. pour $\mathbb R^{\unsure{n|-}}i_*(Q^*)$. Quant à $\deleted{\unsure{D}}\mathbb R^ni_*(\mathbb R j_*(F))$ c'est aussi $\mathbb R^n(ij)_*(F)$, ce qui nous ramène au cas où $U$ est \uline{régulier}. Alors par \uuline{résolution}, on a $i=fi'$, avec $f$ projectif, $X'$ régulier, $i'$ une immersion ouverte à complémentaire un diviseur à croisements normaux

\diagram{\begin{tikzcd}[ampersand replacement=\&,column sep=large,row sep=small] {} \& X'\arrow[d,"f"]\\ U\arrow[ur,hook,"i'"]\arrow[r,hook,no head,"i"'] \& X \end{tikzcd}}

\[ \text{Alors}\qquad R^*fi'_*(F)\Longleftarrow \mathbb R^p f_*(R^q i'_*(F)) \]

et cela nous ramène \inserted{par le lemme 3} au cas où \deleted{\unsure{$\mathbb Q_\ell$} $X$} $X-U$ est un diviseur à croisements normaux. \uuline{Lorsque} $F=\mathbb Q_\ell$, on gagne alors par des calculs explicites \inserted{(th. de \uuline{pureté}).} On va ramener le cas général à celui-ci, en compactifiant $V\deleted{\unsure{/X}}$ en \deleted{\unsure{$X$}}, $g:Y\unsure{\text{---}|\longrightarrow}X$ morphisme projectif (pas lisse a priori), où $F$ est un facteur direct de $\mathbb R^nf_*(\mathbb Q_\ell)$ ; \inserted{on peut supposer dim rel. $g\leq d'$ ;}

\diagram{\begin{tikzcd}[ampersand replacement=\&,column sep=large,row sep=small]
V\arrow[r,hook,no head,"\unsure{d|j}"]\arrow[d,"f"'] \& \widehat Y\arrow[d,"g"]\\
U\arrow[r,hook,no head,"i"'] \& X
\end{tikzcd}}

On peut supposer $F=\mathbb R^n f_*(\mathbb Q_\ell),=\mathbb R^n g_*(\mathbb Q_\ell)|U$. Par des dévissages faciles, utilisant lemme 3 on est ramené au cas où $d'=1$\inserted{ et $f$ à fibres géom. conn.}. Alors la suite spectrale pour $if$ est

% slanted
\margin{\uline{développer}}

% archive sheet number: 134
